% ================================================================
% ==== FILE: content/k_levels/k4.tex
% ================================================================

% ==============================
%  Ontology of Continua — Core
%  K-level Module: K4
%  File: content/k_levels/k4.tex
%  Status: FULLY DEFINED
% ==============================



\paragraph{Canonical tuple projection note.}
Local K-level displays in this module are projections of the canonical public tuple \(K=(\Omega,\partial\Omega,A,\Theta,P,J,C,k,M)\). When a display omits \(M\), reorders components, or uses level-indexed shorthand, the omitted embedding context is inherited from the surrounding level discussion rather than introduced as a competing definition.

\subsubsection{\texorpdfstring{$K_4$}{K_4} Übersicht}
\label{sec:k4-overview}

$K_4$ ist die Ebene, auf der eine \emph{biologische} Struktur zuerst entsteht
möglich. Das entscheidende Ereignis ist das Erscheinen eines \emph{semi-stable
Membran}, die eine neue ontologische Unterscheidung erzeugt:
\[
\text{innen} \;|\; \text{draußen}.
\]

Dies ist eine echte neue Achse des Kontinuums, die Folgendes schafft:
\begin{itemize}
    \item ein begrenztes chemisches Inneres,
    \item selektive Transportregime,
    \item räumlich strukturierte Farbverläufe,
    \item Protohomöostase,
    \item die ersten biologischen Zyklen der Erregung, des Überlebens und der Reparatur.
\end{itemize}

Der Übergang $K_3 \to K_4$ entspricht der Protozellenbildung: einer RAF
Netzwerk, das von einer semipermeablen Membran umgeben ist, deren Stabilität bestimmt wird
durch Schwellenwertbedingungen $\Theta_{\mathrm{mem}}$,
$\Theta_{\mathrm{perm}}$, $\Theta_{\mathrm{osm}}$ und
$\Theta_{\mathrm{curv}}$.

$K_4$ ist das minimale biologische Kontinuum: die einfachste Form der Organisation
und nachhaltiges protozelluläres Leben.

% ================================================================
\subsubsection{Zustandsraum \texorpdfstring{$\Omega(K_4)$}{\Omega(K_4)}}

Der Zustandsraum umfasst:
\[
\Omega(K_4) =
\{(\mathbf{c}_{\mathrm{in}},\mathbf{c}_{\mathrm{out}},
  \partial\Omega,\; \mathbf{g},\; \Gamma_{\mathrm{RAF}},
  \mathbf{p})\},
\]
wo:
\begin{itemize}
    \item $\mathbf{c}_{\mathrm{in}}$: Konzentrationen innerhalb der Membran,
    \item $\mathbf{c}_{\mathrm{out}}$: externe Konzentrationen,
    \item $\partial\Omega$: Membrangrenze mit Patch-Struktur,
    \item $\mathbf{g}$: Gradienten über die Membran (osmotisch, ionisch,
          pH-Wert, Potenzial),
    \item $\Gamma_{\mathrm{RAF}}$: das interne katalytische Netzwerk,
    \item $\mathbf{p}$: Permeabilitäts- und Krümmungsparameter für jeden
          Grenzfleck.
\end{itemize}

Der Zustandsraum ist aufgrund der Membrantopologie höherdimensional als $K_3$
und abschnittsweise Freiheitsgrade.

Zulässige Staaten müssen Folgendes erfüllen:
\begin{itemize}
    \item Membrankohäsion,
    \item kontrollierte Durchlässigkeit,
    \item zerstörungsfreie Farbverläufe,
    \item Lebensfähigkeit des RAF-Netzwerks im Inneren.
\end{itemize}

% ================================================================
\subsubsection{Grenze \texorpdfstring{$\partial\Omega(K_4)$}{\partial\Omega(K_4)}}

Die Membran ist eine \emph{dynamische Grenze}:
\[
\partial\Omega(K_4) = \bigcup_{i=1}^{N_{\mathrm{patch}}}
    \partial\Omega_i.
\]

Jeder Patch $\partial\Omega_i$ hat:
\begin{itemize}
    \item lokale Krümmung $K_i$,
    \item lokale Permeabilität $p_i$,
    \item lokale Spannung $T_i$,
    \item lokale Gradientenbeschränkungen (osmotisch, ionisch, pH),
    \item Phasenzustand (L\textsubscript{α}, L\textsubscript{o},
          L\textsubscript{β}, porös).
\end{itemize}

Ein Patch bricht zusammen, wenn:
\[
T_i > \Theta_{\mathrm{mem},i}
\quad \text{oder} \quad
\Delta \Pi_i > \Theta_{\mathrm{osm},i}
\quad \text{oder} \quad
K_i > \Theta_{\mathrm{curv},i}.
\]

Der Verlust eines kritischen Patches zerstört $\Omega(K_4)$.

% ================================================================
\subsubsection{Achsen \texorpdfstring{$A(K_4)$}{A(K_4)}}

$K_4$ unterstützt mindestens die folgenden Achsen:
\[
A(K_4) = \{ A_{\mathrm{comp}}, A_{\mathrm{boundary}},
           A_{\mathrm{curv}}, A_{\mathrm{perm}}, A_{\mathrm{grad}}\}.
\]

Interpretation:
\begin{itemize}
    \item $A_{\mathrm{comp}}$: Innen-/Außenfachunterscheidung,
    \item $A_{\mathrm{boundary}}$: Freiheitsgrade der Membran,
    \item $A_{\mathrm{curv}}$: Krümmungszustände von Grenzflächen,
    \item $A_{\mathrm{perm}}$: variable Permeabilitätsregime,
    \item $A_{\mathrm{grad}}$: chemische/elektrische/ionische Gradienten.
\end{itemize}

Diese Achsen sind \emph{nicht} auf chemische Achsen von $K_3$ reduzierbar; jeder
stellt neue strukturelle Unterschiede der biologischen Organisation dar.

% ================================================================
\subsubsection{Potentiale \texorpdfstring{$P(K_4)$}{P(K_4)}}

Zu den Potenzialen gehören:
\[
P(K_4) =
\{ \mu_i^{\mathrm{in}},\; \mu_i^{\mathrm{out}},\;
   P_{\mathrm{grad}},\;
   P_{\mathrm{osm}},\;
   P_{\mathrm{curv}},\;
   P_{\mathrm{mem}},\;
   P_{\mathrm{Pumpe}} \}.
\]

Bedeutung:
\begin{itemize}
    \item interne/externe chemische Potentiale,
    \item-Gradientenpotentiale, die den selektiven Transport vorantreiben,
    \item osmotische Druckpotentiale,
    \item Krümmungsenergie von Membranflecken,
    \item Membranspannungspotentiale,
    \item in primitiven Pumpen gespeicherte Energie (falls vorhanden).
\end{itemize}

Diese Potenziale regulieren das Überleben: ein zu großer Gradient oder eine zu große Krümmung
zerstört das System.

% ================================================================
\subsubsection{Schwellenwerte \texorpdfstring{$\Theta(K_4)$}{\Theta(K_4)}}

Schwellenwerte bestimmen die Lebensfähigkeit von Membranen und Protozellen:
\[
\Theta(K_4) =
\{\Theta_{\mathrm{mem}},\;
  \Theta_{\mathrm{perm}},\;
  \Theta_{\mathrm{osm}},\;
  \Theta_{\mathrm{curv}},\;
  \Theta_{\mathrm{pH}},\;
  \Theta_{\mathrm{grad}},\;
  \Theta_{\mathrm{RAF-Überleben}}\}.
\]

Wichtige Interpretationen:
\begin{itemize}
    \item $\Theta_{\mathrm{mem}}$: maximale Membranspannung vorher
          Bruch,
    \item $\Theta_{\mathrm{perm}}$: Grenzen der passiven Permeabilität,
    \item $\Theta_{\mathrm{osm}}$: osmotische Toleranz vor dem Ausbruch,
    \item $\Theta_{\mathrm{curv}}$: Krümmungsinstabilitätsschwelle,
    \item $\Theta_{\mathrm{pH}}$: interner pH-Lebensfähigkeitsbereich,
    \item $\Theta_{\mathrm{grad}}$: Gradienten nachhaltig ohne
          Zusammenbruch,
    \item $\Theta_{\mathrm{RAF-survival}}$: minimale katalytische Aktivität
          für homöostatische Zyklen.
\end{itemize}

Diese Schwellenwerte bestimmen gemeinsam die Grenze lebensfähiger Protozellen
Staaten.

% ================================================================
\subsubsection{Flüsse \texorpdfstring{$J(K_4)$}{J(K_4)}}

Zu den Flüssen in $K_4$ gehören:
\begin{itemize}
    \item Transmembrantransport:
          \[
          J_{\mathrm{ion}},\; J_{\mathrm{Wasser}},\; J_{\mathrm{gelöster Stoff}},
          \]
    \item passive Leckströme,
    \item krümmungsgetriebene Strömungen und Patch-Umordnungen,
    \item interne RAF-reaktive Flüsse,
    \item pumpengetriebene aktive Strömungen (sofern primitive Pumpen vorhanden sind),
    \item Osmotische Schwellungs-Entspannungs-Zyklen.
\end{itemize}

Die Strömungsstabilität erfordert die Kontrolle von:
\[
\Delta V,\; \Delta \Pi,\; \Delta c_i,\; \Updelta \mathrm{pH}.
\]

% ================================================================
\subsubsection{Zyklen \texorpdfstring{$C(K_4)$}{C(K_4)}}

Biologische Zyklen treten zuerst auf:
\[
C(K_4) =
\{ C_{\mathrm{Überleben}},\;
   C_{\mathrm{osmotisch}},\;
   C_{\mathrm{pH-Erholung}},\;
   C_{\mathrm{Leckpumpe}},\;
   C_{\mathrm{RAF-Stabilität}} \}.
\]

Bedeutung:
\begin{itemize}
    \item $C_{\mathrm{survival}}$: Erhaltung des minimalen homöostatischen Zyklus
          Membranintegrität,
    \item $C_{\mathrm{osmotische}}$: wiederholte Schwellungs-Entspannungs-Zyklen,
    \item $C_{\mathrm{pH-recovery}}$: Pufferung und Reaktionsflüsse,
    \item $C_{\mathrm{leak-pump}}$: Gleichgewicht zwischen passiven und aktiven Lecks
          oder semiaktiver Transport,
    \item $C_{\mathrm{RAF-Stabilität}}$: Einspeisung der internen Netzwerkdynamik
          Grenzstabilität.
\end{itemize}

Zyklen definieren die biologische Zeitrichtung und die Protozellpersistenz.

% ================================================================
\subsubsection{Zeit \texorpdfstring{$\tau(K_4)$}{\tau(K_4)}}

Zeit existiert, wenn:
\[
\Pi(C_{\mathrm{Überleben}}) > \Theta_{\mathrm{Zeit}}
\quad\text{und}\quad
\Pi(C_{\mathrm{RAF-Stabilität}}) > \Theta_{\mathrm{Zeit}}.
\]

Also:
\begin{itemize}
    \item biologische Zeit ergibt sich aus Überlebenszyklen,
    \item Die Zeit bricht zusammen, wenn die Membran keinen stabilen Zyklus aufrechterhalten kann
          Betrieb,
    Die zeitliche Reihenfolge von \item ist aufgrund der Grenze stärker als in $K_3$
          Rückkopplungsschleifen.
\end{itemize}

% ================================================================
\subsubsection{Kontinuität \texorpdfstring{$k(K_4)$}{k(K_4)}}

Anwendung der allgemeinen Formel:
\[
k(K_4) =
\frac{\mu(\Omega(K_4))}{\mu(S(K_4))}
\cdot
\frac{|A(K_4)|}{|A(M_4)|}
\cdot
\frac{\sum \mathrm{Stab}(C)}{\sum \mathrm{MaxStab}(C)}
\cdot
(1 - \frac{\sigma(J)}{\sigma_{\max}})
\cdot
(1 - \frac{T}{\Theta_{\mathrm{stab}}})_+.
\]

Implikationen:
\begin{itemize}
    \item Membranstabilität erhöht $k(K_4)$,
    \item Osmotische und Krümmungsschwankungen nehmen $k(K_4)$ ab,
    \item Verlust der Integrität auf Patch-Ebene setzt $k(K_4)=0$,
    \item erfolgreiche Protozellenorganisation entspricht
          $k(K_4)>0$ für längere Zeiträume.
\end{itemize}

% ================================================================
\subsubsection{Strukturelle Spannung \texorpdfstring{$T(K_4)$}{T(K_4)}}

Quellen:
\begin{itemize}
    \item inkompatible osmotische Gradienten,
    \item übermäßige Krümmungsspannung an Grenzflächen,
    \item pH-Stress,
    \item unkontrollierte Durchlässigkeit,
    \item Ungleichgewicht zwischen Leck und Pumpendurchfluss,
    \item unzureichende RAF-Unterstützung.
\end{itemize}

Strukturelles Versagen liegt vor, wenn:
\[
T(K_4) > \Theta_{\mathrm{mem}}.
\]

% ================================================================
\subsubsection{Energie \texorpdfstring{$E(K_4)$}{E(K_4)}}

Die Energiefunktion umfasst:
\[
E(K_4) =
E_{\mathrm{chem}}^{\mathrm{in}} + E_{\mathrm{chem}}^{\mathrm{out}}
+ E_{\mathrm{osm}} + E_{\mathrm{curv}}
+ E_{\mathrm{grad}} + E_{\mathrm{mem}}
+ E_{\mathrm{RAF}}.
\]

Komponenten:
\begin{itemize}
    \item innere und äußere chemische Energien,
    \item osmotische Druckenergie,
    \item Membrankrümmungsenergie,
    \item Gradientenenergie,
    \item Membranspannungsenergie,
    \item RAF-Netzwerk katalytische Energie.
\end{itemize}

$E(K_4)$ bestimmt die Stabilität und Dynamik der Protozelle.

% ================================================================
\subsubsection{Operatoren auf \texorpdfstring{$K_4$}{K_4} (\texorpdfstring{$\Psi$}{\Psi}, \texorpdfstring{$\Phi$}{\Phi}, \texorpdfstring{$\Lambda$}{\Lambda}, \texorpdfstring{$U$}{U}, \texorpdfstring{$\Chi$}{\Chi})}

\begin{itemize}
    \item $\Psi_{4\to5}$ — Entstehung elektrischer Erregbarkeit und
          primitive Ionenkanäle (Übergang zu $K_5$),
    \item $\Phi$ – Entwicklung der Membranbeschränkungen und -zulässigkeit
          Chemie,
    \item $\Lambda$ – strukturelle Kopplung von Membran- und RAF-Sets,
    \item $U$ – Vereinheitlichung von Gradienten- und Membran-Subsystemen,
    \item $\Chi$ – Verzweigung in alternative Protozellentypen (verschiedene
          Membranchemie oder Transportregime).
\end{itemize}

% ================================================================
\subsubsection{Prozesse auf \texorpdfstring{$K_4$}{K_4}}

Charakteristische Prozesse:
\begin{itemize}
    \item Vesikelbildung und -verschluss,
    \item krümmungsbedingte Umordnungen,
    \item osmotische Schwellung und Entspannung,
    \item pH-Pufferzyklen,
    \item Leck-Pumpen-Gleichgewichtsdynamik,
    \item protometabolische Rückkopplungsschleifen,
    \item Entstehung früher Ionenkanäle,
    \item Stabilisierung der elektrischen Achse (vor $K_5$).
\end{itemize}

% ================================================================
\subsubsection{Vorhersagen für \texorpdfstring{$K_4$}{K_4}}

\begin{enumerate}
    \item Membranstabilität erfordert Krümmung und osmotische Schwellenwerte
          in engen Grenzen.
    \item Patch-basierte Spannungsverteilung sagt lokale Flimmermodi voraus.
    \item pH-Erholungszyklen sind für den Fortbestand der Protozelle erforderlich.
    \item RAF-Netzwerke allein können ihre Organisation ohne sie nicht aufrechterhalten
          Grenzstabilität.
    \item Elektrische Gradienten beginnen vor dem vollständigen $K_5$ zu erscheinen.
    \item Der Übergang zu $K_5$ erfolgt erst nach Stabilisierung des Ions
          Kanäle.
\end{enumerate}

% ================================================================
\subsubsection{Experimente für \texorpdfstring{$K_4$}{K_4}}

Empirische Analoga umfassen:
\begin{itemize}
    \item Fettsäurevesikel-Experimente,
    \item Protozell-Osmose-Burst-Tests,
    \item Patch-Level-Fluoreszenzkartierung des Membranflimmerns,
    \item Ionenkanaleinfügung in einfache Vesikel,
    \item In Vesikeln eingeschlossene RAF-Netzwerke,
    \item pH-gepufferte Vesikeldynamik.
\end{itemize}

Diese Tests validieren Schwellenwerte $\Theta_{\mathrm{mem}}$,
$\Theta_{\mathrm{curv}}$, $\Theta_{\mathrm{osm}}$,
$\Theta_{\mathrm{pH}}$.

% ================================================================
\subsubsection{Zusammenbruch und Tod von \texorpdfstring{$K_4$}{K_4}}

Ein Kollaps tritt auf, wenn:
\[
\Omega(K_4) = \varnothing.
\]

Mitwirkende:
\begin{itemize}
    \item Membranbruch,
    \item extremes osmotisches Ungleichgewicht,
    \item krümmungsbedingtes Bersten,
    \item katastrophaler pH-Kollaps,
    \item Leck außer Kontrolle geraten,
    \item Versagen von Überlebenszyklen.
\end{itemize}

Sobald es zusammengebrochen ist, ist ein Übergang zu $K_5$ unmöglich.

% ================================================================
\subsubsection{Falschbarkeit von \texorpdfstring{$K_4$}{K_4}}

Vorhersagen, die experimentell in Frage gestellt werden können:
\begin{enumerate}
    \item die Existenz strenger Membranschwellen für das Überleben,
    \item universelle Form der Osmosekollapskurven,
    \item Notwendigkeit von RAF-unterstützten Pufferzyklen,
    \item Unmöglichkeit einer nachhaltigen Protozellenorganisation ohne
          Krümmungskontrolle auf Patch-Ebene,
    \item frühes Auftreten protoelektrischer Gradienten,
    \item Anforderung eines minimalen Leckage-Pumpen-Gleichgewichts.
\end{enumerate}

Ein Scheitern dieser Vorhersagen würde $K_4$ wie derzeit definiert widerlegen.

% ================================================================
\subsubsection{Verzweigung / ontologische Position von \texorpdfstring{$K_4$}{K_4}}

Verzweigungstypen:
\begin{itemize}
    \item verschiedene Membranchemien (Fettsäure, Phospholipid,
          gemischt),
    \item unterschiedliche Permeabilitätsregime,
    \item alternative RAF-Kerne, die mit der Membrandynamik kompatibel sind,
    \item Divergente protometabolische Strategien.
\end{itemize}

Ontologische Position:
\[
K_3 \to K_4 \to K_5,
\]
mit $K_4$ als erstem biologischen Kontinuum.

% ================================================================
\subsubsection{Beziehung zu M-Räumen}

$M_4$ gibt an:
\begin{itemize}
    \item zulässige Membranmaterialien,
    \item erlaubte Krümmungs- und Durchlässigkeitsbereiche,
    \item Umweltregime (osmotisch, ionisch, pH),
    \item erlaubte Transportvorgänge,
    \item-Einschränkungen für die Bildung von $\partial\Omega$ als stabile Grenze.
\end{itemize}

Die Kompatibilität von $K_4$ mit $M_4$ bestimmt die Lebensfähigkeit der Protozellen.
